\chapter{Conclusion and Future Work}
\label{ch:conclusion}

\section{Summary of Contributions}
\label{sec:contributions}

This project built, benchmarked, and demonstrated an end-to-end neuro-symbolic distillation
pipeline that addresses the token cost problem in enterprise LLM deployments. The system
takes raw supply chain event streams, builds a live knowledge graph, computes structural
deltas, scores each change by business saliency, and produces compact Business Intelligence
Briefings that preserve analytical quality while reducing token usage by up to two orders
of magnitude. The Streamlit dashboard makes the pipeline accessible to non-technical users
through four interactive pages.

The five project objectives from Chapter~\ref{ch:introduction} were met as follows.
The neuro-symbolic pipeline was built and is fully operational, processing synthetic supply
chain events through all five stages with less than 15~ms of overhead before any LLM call.
The 10$\times$ compression target was exceeded at every scale tested, reaching 84.1$\times$
at 1,000 events and 311$\times$ at 3,000 events, while recall remained at 93.8\%---above
the 90\% threshold specified in the objective. The Streamlit dashboard was implemented with
all four pages (Pipeline Demo, Knowledge Graph, Benchmarks, Agent Insights). The benchmark
framework compared the full pipeline against both a naive baseline and an aggregation-only
baseline at five scales. The ground truth evaluation detected all four injected patterns,
with the pipeline uniquely identifying cascade effects that the aggregation baseline missed.

The project's core architectural claim is that symbolic preprocessing of business events
can substitute for the majority of the LLM's in-context reasoning burden. By computing
entity relationships, change magnitudes, and importance scores before the API call, the
pipeline offloads work from the expensive stochastic inference layer to a cheap deterministic
computation layer. The 734-token briefing that an LLM agent receives is not a summarization
of the raw data---it is a structured analytical artifact that already encodes which entities
are affected, how their relationships connect, and how severe the changes are. The agent's
job is to interpret this artifact, not to reconstruct it from scratch.

\section{Limitations}
\label{sec:limitations}

The system as built has four clear limitations that bound the scope of the claims made in
this report.

\textbf{Synthetic data only.} The entire evaluation is on procedurally generated events
using a fixed graph topology and calibrated anomaly injection. Real enterprise event streams
have irregular schemas, missing values, encoding errors, and distribution shifts that the
synthetic generator does not reproduce. The compression ratios and recall figures reported
here represent a best-case scenario relative to what a production deployment would achieve.
How much performance degrades on real data is an open question that future work must address.

\textbf{Single domain.} The knowledge graph schema, saliency formula coefficients, and
distillation format were all designed for supply chain events. Applying the same system to
a different domain---financial transactions, healthcare records, or manufacturing sensor
data---would require re-engineering the graph schema and recalibrating the saliency
weights. The approach generalizes as a concept but is domain-specific as an implementation.

\textbf{LLM API dependency.} The agent layer is a thin wrapper around the Gemini Flash
API. The quality results depend on the specific model's analytical capabilities and are
not reproducible if the model version changes. An open-weights locally-hosted alternative
would give more reproducible and auditable results, at the cost of additional infrastructure.

\textbf{Hand-tuned saliency threshold.} The threshold $\theta = 0.4$ and the coefficient
values $\alpha = 0.35$, $\beta = 0.35$, $\gamma = 0.30$ were tuned by grid search on the
four ground truth patterns. This is a small calibration set. A larger labeled evaluation
set would allow more principled coefficient selection and might yield better recall at
higher compression ratios.

\section{Future Work}
\label{sec:future_work}

The four limitations above point directly to the most valuable directions for future work.

\textbf{Real enterprise data connectors.} The highest-impact extension is replacing the
synthetic data generator with connectors to real ERP and supply chain management systems.
Adapters for SAP, Oracle SCM, or message queue systems like Apache Kafka would let the
pipeline process live event streams and make the compression and recall measurements
meaningful for production deployments. This is also the most straightforward engineering
extension: the pipeline's interface to the data generator is a simple Python generator
function, and replacing it requires no changes to the symbolic or neural layers.

\textbf{Learned saliency via vector embeddings.} The hand-crafted saliency formula is
interpretable but brittle. A learned embedding model that represents event importance in a
continuous space, trained on historical incidents labeled by domain experts, would generalize
better across different supply chain configurations and event distributions. Combining the
structural graph signals with learned semantic similarity---a hybrid symbolic-neural score---is
a natural next step that the existing architecture can accommodate without restructuring.

\textbf{Multi-domain graph schemas.} Generalizing the graph schema to support additional
domains (financial transaction networks, healthcare supply chains, manufacturing bills of
materials) would test whether the architectural pattern transfers and would increase the
system's utility as a reusable open-source tool. Domain-specific schema adapters could
be implemented as plugins without modifying the core pipeline.

\textbf{On-device distillation for edge deployment.} The pipeline currently runs on a
server and calls a cloud LLM API. For deployments where data privacy or network availability
is a constraint, running the full pipeline---including a small locally-hosted language
model---on edge hardware would make the system viable in air-gapped environments. The
pipeline's sub-15~ms processing overhead suggests that the bottleneck in such a deployment
would be the local model inference speed, not the symbolic stages.
